\chapter{Additional Result Visualizations}
\label{app:additional-result-visualizations}

This appendix provides distributional and observation-level views that complement the principal result figures. Each figure states whether its unit of analysis is a question, a question--concept observation, or a concept aggregated across questions.

\section{Causal-Effect Shifts under Joint Interventions}
\label{app:ce-shift-rainclouds}

Figure~\ref{fig:appendix-ce-shift-rainclouds} complements the faithfulness distributions in Figure~\ref{fig:question-level-faithfulness-rainclouds}. For each question, the plotted value is the arithmetic mean of the concept-level CE values. The atomic and interaction-aware values originate from the same joint FELICS run: the former averages \texttt{atomic\_kl\_div}, whereas the latter averages \texttt{adjusted\_shapley\_kl\_div}. Thus, concepts, responses, and explanation-implied effects are held fixed within each comparison.

\ifthesislocalbuild
  \definecolor{CeAtomic}{HTML}{2878B5}
\definecolor{CeJoint}{HTML}{E37A26}

% tikzviolinplots computes each KDE in \violin@kde. Preserve the KDE before
% the second violin overwrites it, following the summary raincloud figure.
\newcount\cekderow
\makeatletter
\newcommand{\CeCaptureCurrentKDE}[1]{%
  \pgfplotstablegetrowsof{\violin@kde}%
  \pgfmathtruncatemacro{\cekdelast}{\pgfplotsretval-1}%
  \global\def#1{list kdecol\\}%
  \cekderow=0
  \loop\ifnum\cekderow<\numexpr\cekdelast+1\relax
    \pgfplotstablegetelem{\the\cekderow}{list}\of\violin@kde
    \edef\cekdevalue{\pgfplotsretval}%
    \pgfplotstablegetelem{\the\cekderow}{kdecol}\of\violin@kde
    \xdef#1{#1\cekdevalue\space\pgfplotsretval\\}%
    \advance\cekderow by 1
  \repeat
}
\newcommand{\CeMaterializeKDE}[2]{%
  \begingroup
  \edef\celoadtable{%
    \endgroup
    \noexpand\pgfplotstableread[row sep=\noexpand\\]%
      {\expandafter\unexpanded\expandafter{#1}}\noexpand#2}%
  \celoadtable
}
\makeatother

% #1 panel title, #2 CSV, #3 file-safe id, #4 x maximum, #5 x ticks,
% #6 tick precision. The two models use the same scale within each dataset.
\newcommand{\CeRaincloudPanel}[6]{%
  \begin{minipage}[t]{.495\textwidth}
  \centering
  {\small\bfseries #1\par}
  \vspace{.35em}
  \edef\CeDataHash{\pdfmdfivesum file {#2}}%
  \tikzsetnextfilename{ce-shift-#3-\CeDataHash}%
  \begin{tikzpicture}
    \begin{scope}
    \clip (0,0) rectangle (0,0);
    \violinsetoptions[scaled,reverseaxis]{
      width=8.55cm,
      height=5.40cm,
      overlay,
      xmin=0,
      xmax=#4,
      ymin=.10,
      ymax=7.25,
      axis lines=none,
    }
    \violinplot[col sep=comma,index=atomic_ce,relative position=5.05,
      color=white,label={},samples=30]{#2}
    \CeCaptureCurrentKDE{\AtomicCeKDE}
    \violinplot[col sep=comma,index=interaction_ce,relative position=1.65,
      color=white,label={},samples=30]{#2}
    \CeCaptureCurrentKDE{\JointCeKDE}
    \end{scope}
    \CeMaterializeKDE{\AtomicCeKDE}{\AtomicCeKDETable}
    \CeMaterializeKDE{\JointCeKDE}{\JointCeKDETable}

    \pgfresetboundingbox
    \begin{axis}[
      width=8.55cm,
      height=5.55cm,
      xmin=0,
      xmax=#4,
      ymin=.10,
      ymax=7.25,
      axis x line*=bottom,
      axis y line=none,
      xtick={#5},
      scaled ticks=false,
      tick label style={/pgf/number format/fixed,
        /pgf/number format/precision=#6},
      xmajorgrids=true,
      grid style={gray!18},
      clip=true,
    ]
      \addplot[draw=CeAtomic,fill=CeAtomic,fill opacity=.48]
        table[x=list,y expr=5.15+1.6*(\thisrow{kdecol}-5.05)]
          {\AtomicCeKDETable};
      \addplot[only marks,mark=*,mark size=1.15pt,
        draw=CeAtomic,fill=CeAtomic,fill opacity=.62,
        y filter/.expression={4.35+.16*(mod(\coordindex*7,11)/10-.5)}]
        table[col sep=comma,x=atomic_ce,y expr=4.35]{#2};
      \addplot+[boxplot={draw direction=x,draw position=4.02},
        mark=none,draw=CeAtomic,fill=CeAtomic,fill opacity=.32,
        thick,boxplot/box extend=.18]
        table[col sep=comma,y=atomic_ce]{#2};
      \addplot[draw=CeJoint,fill=CeJoint,fill opacity=.48]
        table[x=list,y expr=1.75+1.6*(\thisrow{kdecol}-1.65)]
          {\JointCeKDETable};
      \addplot[only marks,mark=*,mark size=1.15pt,
        draw=CeJoint,fill=CeJoint,fill opacity=.66,
        y filter/.expression={.92+.16*(mod(\coordindex*7,11)/10-.5)}]
        table[col sep=comma,x=interaction_ce,y expr=.92]{#2};
      \addplot+[boxplot={draw direction=x,draw position=.55},
        mark=none,draw=CeJoint,fill=CeJoint,fill opacity=.34,
        thick,solid,boxplot/box extend=.18]
        table[col sep=comma,y=interaction_ce]{#2};
      \node[anchor=west,font=\scriptsize\bfseries,text=CeAtomic]
        at (rel axis cs:.03,.937) {Atomic};
      \node[anchor=west,font=\scriptsize\bfseries,text=CeJoint]
        at (rel axis cs:.03,.462) {Joint};
    \end{axis}
  \end{tikzpicture}%
  \end{minipage}%
}

\tikzexternalenable
\begin{figure}[p]
  \centering
  \CeRaincloudPanel{BBQ --- GPT-OSS 120B}%
    {data/appendix/ce_shifts/bbq_gptoss.csv}{bbq-gptoss}{.6}{0,.2,.4,.6}{1}%
  \hfill
  \CeRaincloudPanel{BBQ --- Mistral Medium 3.5}%
    {data/appendix/ce_shifts/bbq_mistral.csv}{bbq-mistral}{.6}{0,.2,.4,.6}{1}%

  \par\vspace{1.2em}
  \CeRaincloudPanel{MedQA --- GPT-OSS 120B}%
    {data/appendix/ce_shifts/medqa_gptoss.csv}{medqa-gptoss}{.06}{0,.02,.04,.06}{2}%
  \hfill
  \CeRaincloudPanel{MedQA --- Mistral Medium 3.5}%
    {data/appendix/ce_shifts/medqa_mistral.csv}{medqa-mistral}{.06}{0,.02,.04,.06}{2}%

  \par\vspace{1.2em}
  \CeRaincloudPanel{New German Credit --- GPT-OSS 120B}%
    {data/appendix/ce_shifts/new_german_credit_gptoss.csv}{credit-gptoss}{1.5}{0,.5,1,1.5}{1}%
  \hfill
  \CeRaincloudPanel{New German Credit --- Mistral Medium 3.5}%
    {data/appendix/ce_shifts/new_german_credit_mistral.csv}{credit-mistral}{1.5}{0,.5,1,1.5}{1}%

  \caption[Raincloud distributions of atomic and interaction-aware CE]{Distribution of question-averaged causal effects obtained from the atomic and interaction-aware CE vectors of the same joint FELICS runs. Axis scales are shared between models within each dataset but differ across datasets because their CE magnitudes differ substantially.}
  \label{fig:appendix-ce-shift-rainclouds}
\end{figure}
\tikzexternaldisable

\else
  \typeout{Non-local build detected: skipping graphs/results/appendix_ce_rainclouds.tex}
  \begin{tcolorbox}[
    colback=orange!12,
    colframe=orange!80!black,
    fonttitle=\bfseries,
    title={Compilation warning}
  ]
    This part of the document cannot be compiled on Overleaf and has therefore
    been omitted from this build.
  \end{tcolorbox}
\fi
\FloatBarrier

\section{CE Agreement between Implementations}
\label{app:rq2-ce-identity}

Figures~\ref{fig:appendix-rq2-ce-natural-language} and~\ref{fig:appendix-rq2-ce-credit} complement the faithfulness identity plots used to answer Research Question~\ref{rq:convergence}. Each point corresponds to a question available in both independently sampled runs. Its coordinates are the question-level means used by the CE equivalence and paired-difference tests: the mean \texttt{kl\_div} from the Matton implementation and the mean \texttt{adjusted\_shapley\_kl\_div} from FELICS with $n_{\max}=1$.

% Appendix counterpart to the faithfulness identity plots in RQ2.
% Reusable two-dimensional result plot.
% Callers provide data and presentation details through named keys so that
% result sections only need to define their local conventions.
\definecolor{ResultScatterPoints}{HTML}{2878B5}
\definecolor{ResultScatterLine}{HTML}{E37A26}

\newif\ifResultScatterScaleOnlyAxis

\pgfkeys{
  /result scatter/.is family,
  /result scatter/.cd,
  title/.store in=\ResultScatterTitle,
  data/.store in=\ResultScatterData,
  x column/.store in=\ResultScatterXColumn,
  y column/.store in=\ResultScatterYColumn,
  x label/.store in=\ResultScatterXLabel,
  y label/.store in=\ResultScatterYLabel,
  text below/.store in=\ResultScatterTextBelow,
  xmin/.store in=\ResultScatterXMin,
  xmax/.store in=\ResultScatterXMax,
  ymin/.store in=\ResultScatterYMin,
  ymax/.store in=\ResultScatterYMax,
  line slope/.store in=\ResultScatterLineSlope,
  line intercept/.store in=\ResultScatterLineIntercept,
  line style/.store in=\ResultScatterLineStyle,
  minipage width/.store in=\ResultScatterMinipageWidth,
  plot width/.store in=\ResultScatterPlotWidth,
  plot height/.store in=\ResultScatterPlotHeight,
  mark size/.store in=\ResultScatterMarkSize,
  below skip/.store in=\ResultScatterBelowSkip,
  title gap/.store in=\ResultScatterTitleGap,
  x tick precision/.store in=\ResultScatterXTickPrecision,
  y tick precision/.store in=\ResultScatterYTickPrecision,
  scale only axis/.is if=ResultScatterScaleOnlyAxis,
}

\newcommand{\ResultScatterPlot}[1]{%
  \begingroup
  \pgfkeys{/result scatter/.cd,
    title={},
    data={},
    x column={},
    y column={},
    x label={},
    y label={},
    text below={},
    xmin=0,
    xmax=1,
    ymin=0,
    ymax=1,
    line slope=1,
    line intercept=0,
    line style=solid,
    minipage width=0.485\textwidth,
    plot width=\linewidth,
    plot height=5.45cm,
    mark size=1.65pt,
    below skip=-0.25em,
    title gap=2mm,
    x tick precision=3,
    y tick precision=3,
    scale only axis=false,#1
  }%
  \begin{minipage}[t]{\ResultScatterMinipageWidth}
    \centering
    \ifResultScatterScaleOnlyAxis
      \pgfplotsset{result scatter sizing/.style={scale only axis}}
    \else
      \pgfplotsset{result scatter sizing/.style={}}
    \fi
    \begin{tikzpicture}
      \begin{axis}[
        width=\ResultScatterPlotWidth,
        height=\ResultScatterPlotHeight,
        result scatter sizing,
        title={\ResultScatterTitle},
        title style={
          at={(axis description cs:0.5,1)},
          anchor=south,
          yshift=\ResultScatterTitleGap,
          font=\small\bfseries,
        },
        xmin=\ResultScatterXMin,
        xmax=\ResultScatterXMax,
        ymin=\ResultScatterYMin,
        ymax=\ResultScatterYMax,
        xlabel={\ResultScatterXLabel},
        ylabel={\ResultScatterYLabel},
        scaled x ticks=false,
        scaled y ticks=false,
        xticklabel style={
          font=\scriptsize,
          /pgf/number format/fixed,
          /pgf/number format/precision=\ResultScatterXTickPrecision,
        },
        yticklabel style={
          font=\scriptsize,
          /pgf/number format/fixed,
          /pgf/number format/precision=\ResultScatterYTickPrecision,
        },
        label style={font=\scriptsize},
        grid=major,
        grid style={gray!18},
        axis line style={gray!60},
        clip=true,
      ]
        \addplot[
          only marks,
          mark=*,
          mark size=\ResultScatterMarkSize,
          draw=ResultScatterPoints,
          fill=ResultScatterPoints,
          fill opacity=0.68,
        ] table[
          col sep=comma,
          x=\ResultScatterXColumn,
          y=\ResultScatterYColumn
        ] {\ResultScatterData};
        \addplot[
          ResultScatterLine,
          very thick,
          \ResultScatterLineStyle,
          domain=\ResultScatterXMin:\ResultScatterXMax,
          samples=2
        ] {\ResultScatterLineIntercept + \ResultScatterLineSlope*x};
      \end{axis}
    \end{tikzpicture}
    \vspace{\ResultScatterBelowSkip}
    {\scriptsize \ResultScatterTextBelow\par}
  \end{minipage}%
  \endgroup
}


\newcommand{\HypTwoCeIdentityPlot}[2]{%
  \ResultScatterPlot{x column=matton_ce,
    y column=felics_ce,
    x label={$\overline{\mathrm{CE}}_{\mathrm{Matton}}(\mathbf{x})$},
    y label={$\overline{\mathrm{CE}}_{\mathrm{FELICS},\,n_{\max}=1}(\mathbf{x})$},
    line slope=1,
    line intercept=0,
    line style=dashed,
    minipage width=0.47\textwidth,
    plot width=5.55cm,
    plot height=5.55cm,
    mark size=1.8pt,
    below skip=-0.2em,
    scale only axis=true,
    text below={Shared questions: $n=#1$},
    #2
  }%
}

\begin{figure}[p]
  \centering
  \HypTwoCeIdentityPlot{28}{title={BBQ --- GPT-OSS 120B},
    data=data/rq2/bbq_gptoss_ce.csv,
    xmin=0, xmax=.8,
    ymin=0, ymax=.8
  }
  \hfill
  \HypTwoCeIdentityPlot{30}{title={BBQ --- Mistral Medium 3.5},
    data=data/rq2/bbq_mistral_ce.csv,
    xmin=0, xmax=.8,
    ymin=0, ymax=.8
  }

  \par\vspace{1.1em}
  \HypTwoCeIdentityPlot{29}{title={MedQA --- GPT-OSS 120B},
    data=data/rq2/medqa_gptoss_ce.csv,
    xmin=0, xmax=.25,
    ymin=0, ymax=.25
  }
  \hfill
  \HypTwoCeIdentityPlot{24}{title={MedQA --- Mistral Medium 3.5},
    data=data/rq2/medqa_mistral_ce.csv,
    xmin=0, xmax=.25,
    ymin=0, ymax=.25
  }

  \caption[CE identity plots for the natural-language RQ2 comparisons]{Question-averaged causal effects produced by FELICS with $n_{\max}=1$ and the original Matton implementation for the natural-language comparisons. The average is taken across the concepts retained for each question in the respective run.}
  \label{fig:appendix-rq2-ce-natural-language}
\end{figure}

\begin{figure}[p]
  \centering
  \HypTwoCeIdentityPlot{30}{title={New German Credit --- GPT-OSS 120B},
    data=data/rq2/new_german_credit_gptoss_ce.csv,
    xmin=0, xmax=1.1,
    ymin=0, ymax=1.1
  }
  \hfill
  \HypTwoCeIdentityPlot{30}{title={New German Credit --- Mistral Medium 3.5},
    data=data/rq2/new_german_credit_mistral_ce.csv,
    xmin=0, xmax=1.1,
    ymin=0, ymax=1.1
  }

  \caption[CE identity plots for the New German Credit RQ2 comparisons]{Question-averaged causal effects produced by FELICS with $n_{\max}=1$ and the tabular adaptation of the Matton implementation for New German Credit.}
  \label{fig:appendix-rq2-ce-credit}
\end{figure}

\FloatBarrier

\section{Question-Level CE Agreement under Joint Attribution}
\label{app:ce-shift-identity}

Figures~\ref{fig:appendix-ce-shift-identity-natural-language} and~\ref{fig:appendix-ce-shift-identity-credit} provide a paired robustness view of the CE distributions in Figure~\ref{fig:appendix-ce-shift-rainclouds}. Each point represents one question from the same joint FELICS run. Its coordinates are the question-level means of \texttt{atomic\_kl\_div} and \texttt{adjusted\_shapley\_kl\_div}; marker area increases with the number of concepts contributing to both means. The identity line therefore shows whether questions with high atomic CE also retain high interaction-aware CE.

The upper-tail questions generally remain comparatively high under both attribution methods. The marker sizes do not reveal a consistent tendency for questions with more concepts to have larger mean CE; for New German Credit, the concept count is fixed and cannot account for the observed variation.

\definecolor{CeShiftIdentityPoints}{HTML}{2878B5}
\definecolor{CeShiftIdentityLine}{HTML}{E37A26}

% #1 title, #2 data, #3 minimum, #4 maximum, #5 annotation.
\newcommand{\CeShiftIdentityPlot}[5]{%
  \begin{minipage}[t]{0.47\textwidth}
    \centering
    {\small\bfseries #1\par}
    \vspace{0.1em}
    \begin{tikzpicture}
      \begin{axis}[
        width=5.55cm,
        height=5.55cm,
        scale only axis,
        xmin=#3,
        xmax=#4,
        ymin=#3,
        ymax=#4,
        xlabel={$\overline{\mathrm{CE}}_{\mathrm{atomic}}(\mathbf{x})$},
        ylabel={$\overline{\mathrm{CE}}_{\mathrm{joint}}(\mathbf{x})$},
        scaled x ticks=false,
        scaled y ticks=false,
        tick label style={font=\scriptsize},
        label style={font=\scriptsize},
        grid=major,
        grid style={gray!18},
        axis line style={gray!60},
        clip=true,
      ]
        \addplot[
          scatter,
          only marks,
          mark=*,
          draw=CeShiftIdentityPoints,
          fill=CeShiftIdentityPoints,
          fill opacity=0.62,
          scatter/use mapped color={
            draw=CeShiftIdentityPoints,
            fill=CeShiftIdentityPoints
          },
          visualization depends on={\thisrow{marker_size}\as\perpointmarksize},
          scatter/@pre marker code/.append style={/tikz/mark size=\perpointmarksize pt},
        ] table[
          col sep=comma,
          x=atomic_ce,
          y=interaction_ce,
          meta=concept_count
        ] {#2};
        \addplot[
          CeShiftIdentityLine,
          very thick,
          dashed,
          domain=#3:#4,
          samples=2
        ] {x};
      \end{axis}
    \end{tikzpicture}
    \vspace{-0.2em}
    {\scriptsize #5\par}
  \end{minipage}%
}

\newcommand{\CeShiftConceptCountLegend}{%
  \begin{tikzpicture}[baseline=-0.6ex]
    \node[anchor=east,font=\scriptsize] at (-0.15,0) {Concept count:};
    \foreach \x/\size/\label in {0/1.35/4,1.15/1.86/12,2.45/2.60/31}{
      \draw[CeShiftIdentityPoints,fill=CeShiftIdentityPoints,fill opacity=0.62]
        (\x,0) circle[radius=\size pt];
      \node[anchor=west,font=\scriptsize] at (\x+0.10,0) {$\label$};
    }
  \end{tikzpicture}%
}

\begin{figure}[p]
  \centering
  \CeShiftIdentityPlot{BBQ --- GPT-OSS 120B}%
    {data/appendix/ce_shift_identity/bbq_gptoss.csv}{0}{0.6}%
    {$n=30$ questions; $k=4$--$6$ concepts}
  \hfill
  \CeShiftIdentityPlot{BBQ --- Mistral Medium 3.5}%
    {data/appendix/ce_shift_identity/bbq_mistral.csv}{0}{0.6}%
    {$n=30$ questions; $k=4$--$6$ concepts}

  \par\vspace{1.1em}
  \CeShiftIdentityPlot{MedQA --- GPT-OSS 120B}%
    {data/appendix/ce_shift_identity/medqa_gptoss.csv}{0}{0.06}%
    {$n=29$ questions; $k=8$--$31$ concepts}
  \hfill
  \CeShiftIdentityPlot{MedQA --- Mistral Medium 3.5}%
    {data/appendix/ce_shift_identity/medqa_mistral.csv}{0}{0.06}%
    {$n=24$ questions; $k=6$--$16$ concepts}

  \par\vspace{0.7em}
  \CeShiftConceptCountLegend
  \caption[Question-level CE agreement under atomic and joint attribution]{Question-level mean causal effects under atomic and interaction-aware attribution for BBQ and MedQA. The dashed line denotes equality. Marker area increases with the number of concepts contributing to the question-level mean.}
  \label{fig:appendix-ce-shift-identity-natural-language}
\end{figure}

\begin{figure}[p]
  \centering
  \CeShiftIdentityPlot{New German Credit --- GPT-OSS 120B}%
    {data/appendix/ce_shift_identity/new_german_credit_gptoss.csv}{0}{1.5}%
    {$n=30$ questions; $k=12$ concepts}
  \hfill
  \CeShiftIdentityPlot{New German Credit --- Mistral Medium 3.5}%
    {data/appendix/ce_shift_identity/new_german_credit_mistral.csv}{0}{1.5}%
    {$n=30$ questions; $k=12$ concepts}

  \caption[Question-level CE agreement for New German Credit]{Question-level mean causal effects under atomic and interaction-aware attribution for New German Credit. Every question contains the same 12 concepts, so marker size is constant.}
  \label{fig:appendix-ce-shift-identity-credit}
\end{figure}

\FloatBarrier

\section{Concept-Level Effects of Shapley Normalization}
\label{app:rq3-normalization-identity}

Figures~\ref{fig:appendix-rq3-normalization-natural-language} and~\ref{fig:appendix-rq3-normalization-credit} complement the rank-correlation analysis for Research Question~\ref{rq:ranking}. Each point is one concept in one question, with the raw Shapley CE on the horizontal axis and the adjusted Shapley CE after normalization on the vertical axis. These identity plots visualize changes in CE magnitude rather than constituting additional rank-preservation tests.

Figure~\ref{fig:appendix-rq3-credit-concept-delta} summarizes the New German Credit adjustments by concept. For each of the 12 shared concepts, $\Delta\mathrm{CE}=\mathrm{CE}_{\mathrm{adjusted}}-\mathrm{CE}_{\mathrm{raw}}$ is averaged across the 30 evaluated questions. This aggregation identifies concepts that are affected disproportionately on average while retaining the observation-level view in Figure~\ref{fig:appendix-rq3-normalization-credit}.

For GPT-OSS 120B, the largest mean increases occur for Household Income, Liquid Assets, and SCHUFA Score. For Mistral Medium 3.5, the positive adjustment is distributed more broadly across the concepts, while Household Size remains unchanged on average.

\definecolor{NormalizationIdentityPoints}{HTML}{2878B5}
\definecolor{NormalizationIdentityLine}{HTML}{E37A26}

% #1 title, #2 data, #3 minimum, #4 maximum, #5 observations, #6 rho.
\newcommand{\NormalizationIdentityPlot}[6]{%
  \begin{minipage}[t]{0.47\textwidth}
    \centering
    {\small\bfseries #1\par}
    \vspace{0.1em}
    \begin{tikzpicture}
      \begin{axis}[
        width=5.55cm,
        height=5.55cm,
        scale only axis,
        xmin=#3,
        xmax=#4,
        ymin=#3,
        ymax=#4,
        xlabel={$\mathrm{CE}_{\mathrm{raw}}$},
        ylabel={$\mathrm{CE}_{\mathrm{adjusted}}$},
        scaled x ticks=false,
        scaled y ticks=false,
        tick label style={font=\scriptsize},
        label style={font=\scriptsize},
        grid=major,
        grid style={gray!18},
        axis line style={gray!60},
        clip=true,
      ]
        \addplot[
          only marks,
          mark=*,
          mark size=0.95pt,
          draw=NormalizationIdentityPoints,
          draw opacity=0.36,
          fill=NormalizationIdentityPoints,
          fill opacity=0.30,
        ] table[
          col sep=comma,
          x=raw_ce,
          y=adjusted_ce
        ] {#2};
        \addplot[
          NormalizationIdentityLine,
          very thick,
          dashed,
          domain=#3:#4,
          samples=2
        ] {x};
      \end{axis}
    \end{tikzpicture}
    \vspace{-0.2em}
    {\scriptsize Concept observations: $n=#5$; Spearman $\rho=#6$\par}
  \end{minipage}%
}

\begin{figure}[p]
  \centering
  \NormalizationIdentityPlot{BBQ --- GPT-OSS 120B}%
    {data/appendix/rq3_normalization_identity/bbq_gptoss.csv}{-0.3}{1.1}{137}{0.937}
  \hfill
  \NormalizationIdentityPlot{BBQ --- Mistral Medium 3.5}%
    {data/appendix/rq3_normalization_identity/bbq_mistral.csv}{-0.3}{1.1}{137}{0.903}

  \par\vspace{1.1em}
  \NormalizationIdentityPlot{MedQA --- GPT-OSS 120B}%
    {data/appendix/rq3_normalization_identity/medqa_gptoss.csv}{-0.05}{0.6}{435}{0.938}
  \hfill
  \NormalizationIdentityPlot{MedQA --- Mistral Medium 3.5}%
    {data/appendix/rq3_normalization_identity/medqa_mistral.csv}{-0.05}{0.6}{257}{0.963}

  \caption[Concept-level effects before and after normalization]{Concept-level raw and adjusted Shapley causal effects for BBQ and MedQA. Each translucent point represents one concept in one question. The dashed line denotes unchanged CE magnitude; darker regions indicate overlapping observations.}
  \label{fig:appendix-rq3-normalization-natural-language}
\end{figure}

\begin{figure}[p]
  \centering
  \NormalizationIdentityPlot{New German Credit --- GPT-OSS 120B}%
    {data/appendix/rq3_normalization_identity/new_german_credit_gptoss.csv}{-1}{2.7}{360}{0.875}
  \hfill
  \NormalizationIdentityPlot{New German Credit --- Mistral Medium 3.5}%
    {data/appendix/rq3_normalization_identity/new_german_credit_mistral.csv}{-1}{2.7}{360}{0.821}

  \caption[Concept-level normalization effects for New German Credit]{Concept-level raw and adjusted Shapley causal effects for New German Credit. Broader displacement from the identity line complements the lower rank correlations reported for Research Question~\ref{rq:ranking}.}
  \label{fig:appendix-rq3-normalization-credit}
\end{figure}

\FloatBarrier
\definecolor{CreditDeltaBars}{HTML}{2878B5}

% #1 title, #2 data.
\newcommand{\CreditConceptDeltaPlot}[2]{%
  {\small\bfseries #1\par}
  \vspace{0.2em}
  \begin{tikzpicture}
    \begin{axis}[
      width=0.77\textwidth,
      height=6.7cm,
      xbar,
      bar width=6pt,
      xmin=0,
      xmax=0.4,
      xtick={0,0.1,0.2,0.3,0.4},
      ymin=0.4,
      ymax=12.6,
      y dir=reverse,
      ytick={1,...,12},
      yticklabels={
        Credit Purpose,
        Credit Amount,
        Credit Term (months),
        Liquid Assets,
        Household Income,
        SCHUFA Score (\%),
        Real Estate Value,
        Applicant Age,
        Lives In Own House,
        Current Monthly Debt Burden,
        Job Category,
        Household Size
      },
      xlabel={$\overline{\Delta\mathrm{CE}}=\overline{\mathrm{CE}_{\mathrm{adjusted}}-\mathrm{CE}_{\mathrm{raw}}}$},
      scaled x ticks=false,
      xticklabel style={font=\scriptsize,/pgf/number format/fixed,/pgf/number format/precision=1},
      yticklabel style={font=\scriptsize},
      label style={font=\scriptsize},
      xmajorgrids,
      grid style={gray!18},
      axis line style={gray!60},
      axis y line*=left,
      axis x line*=bottom,
      clip=false,
    ]
      \addplot[
        CreditDeltaBars,
        fill=CreditDeltaBars,
        fill opacity=0.72,
        draw=CreditDeltaBars,
      ] table[
        col sep=comma,
        x=mean_delta_ce,
        y=concept_id
      ] {#2};
    \end{axis}
  \end{tikzpicture}%
}

\begin{figure}[p]
  \centering
  \CreditConceptDeltaPlot{GPT-OSS 120B}%
    {data/appendix/rq3_credit_concept_delta/new_german_credit_gptoss.csv}

  \par\vspace{1.0em}
  \CreditConceptDeltaPlot{Mistral Medium 3.5}%
    {data/appendix/rq3_credit_concept_delta/new_german_credit_mistral.csv}

  \caption[Mean normalization adjustment by New German Credit concept]{Mean normalization adjustment for each New German Credit concept, averaged across the 30 questions. Larger bars identify concepts whose assigned causal effect increases more strongly under the adjusted Shapley normalization.}
  \label{fig:appendix-rq3-credit-concept-delta}
\end{figure}

